\documentclass{report}
\usepackage[utf8]{inputenc}
\usepackage[spanish, es-noshorthands]{babel}
\usepackage{amsmath}
\usepackage{amssymb}
\usepackage{amsthm}
\usepackage{float}
\usepackage{tikz}
\usetikzlibrary{automata, positioning}
\usepackage{clrscode3e}

\theoremstyle{definition}
\newtheorem{definicion}{Definición}[chapter]
\theoremstyle{plain}
\newtheorem{teorema}{Teorema}[chapter]
\newtheorem{lema}{Lema}[chapter]
\theoremstyle{remark}
\newtheorem{ejemplo}{Ejemplo}[chapter]

\begin{document}

\setcounter{chapter}{3}
\chapter{Construcción de la gramática producto de una gramática de concatenación de rango y un lenguaje regular}
\label{cap:interseccion}

En el capítulo~\ref{cap:rcg-sat} se presentaron las gramáticas de
concatenación de rango y se revisaron los antecedentes que abordan el
SAT mediante formalismos de la teoría de lenguajes.
En~\cite{fernandezarias2007, aguileralopez2016} se asocia a cada
fórmula booleana un autómata finito que reconoce sus interpretaciones
verdaderas y una gramática que garantiza que las instancias de una
misma variable tomen el mismo valor, y la satisfacibilidad se decide
mediante el problema del vacío sobre la intersección de ambos
lenguajes. Llevar esta estrategia a las gramáticas de concatenación de
rango exige construir una gramática de concatenación de rango que
reconozca la intersección del lenguaje de otra con un lenguaje
regular.

La literatura provee tales construcciones para las
gramáticas libres del contexto~\cite{barhillel1961} y las de
concatenación de rango simple~\cite{bertsch2001rcg}. En las gramáticas
de concatenación de rango, una misma variable puede aparecer en varias
posiciones de una cláusula, lo que introduce una dificultad que las
construcciones anteriores no contemplan. Este capítulo aporta una
construcción que, a partir de una gramática de concatenación de rango
$G$ y un autómata finito $A$, produce una nueva gramática de
concatenación de rango $G'$, la \emph{gramática
producto}\footnote{El término se usa por analogía con el autómata
producto, la construcción estándar para la intersección de dos
lenguajes regulares~\cite{hopcroft2006introduction}.},
cuyo lenguaje es exactamente $L(G) \cap L(A)$.

\section{El rango sobre un autómata}

Antes de construir la gramática producto conviene precisar qué
cadenas pertenecen a $L(G) \cap L(A)$ y qué exige reconocerlas. La
intersección contiene las cadenas que la gramática reconoce y el
autómata acepta simultáneamente. Una cadena
queda fuera en cuanto uno de los dos deja de reconocerla, de modo que
decidir si pertenece a la intersección obliga a comprobar las dos
condiciones sobre la misma cadena.

Comprobar cada condición por separado es directo. Sea la gramática
$G_0 = (N, T, V, P, Par)$, donde $N = \{Par\}$, $T = \{a\}$, $V = \{X\}$
y las cláusulas de $P$ son
\[
  Par(aaX) \to Par(X) \quad\text{y}\quad Par(\varepsilon) \to \varepsilon.
\]
$G_0$ genera el lenguaje $\{a^{2n} \mid n \ge 0\}$, esto es, las cadenas
formadas por $a$ de longitud par.

El autómata $A_0$ sobre el alfabeto $\{a\}$ tiene estados
$\{q_0, q_1, q_2\}$, transiciones $\delta(q_0, a) = q_1$,
$\delta(q_1, a) = q_2$ y $\delta(q_2, a) = q_0$, estado inicial $q_0$ y
estado final $q_0$, y reconoce las cadenas cuya longitud es múltiplo de
tres.

\begin{center}
\begin{tikzpicture}[->, >=stealth, shorten >=1pt, auto, node distance=2.8cm, semithick]
  \node[state, initial, accepting] (q0) {$q_0$};
  \node[state] (q1) [right=of q0] {$q_1$};
  \node[state] (q2) [right=of q1] {$q_2$};
  \path (q0) edge              node {$a$} (q1)
        (q1) edge              node {$a$} (q2)
        (q2) edge [bend left=40] node {$a$} (q0);
\end{tikzpicture}
\end{center}

Ante la cadena $aaaaaa$, el autómata la acepta porque, al leer sus
símbolos, termina en $q_0$, y la gramática la reconoce porque, al aplicar
sus cláusulas, la derivación termina en la cadena vacía.

La dificultad surge al exigir las dos condiciones a la vez sobre la
misma cadena. El reconocimiento en el autómata avanza estado por estado;
el de la gramática divide la cadena en rangos. Para coordinarlos haría
falta saber, de cada rango, en qué estado está el autómata al empezar y
al terminar de leer la subcadena que el rango delimita. El rango $(i, j)$
no contiene esa información: marca posiciones dentro de una cadena fija,
y no registra los estados del autómata.

Hace falta entonces una noción de rango que cargue la información del
autómata. Para ello, cada subcadena se caracteriza por el par de estados
entre los que el autómata transita al leerla. Para fijar ese par hace
falta saber en qué estado termina el autómata al comenzar en un estado y
leer una cadena, lo que da lugar a la siguiente definición.

\begin{definicion}
Sea $A = (Q, \Sigma, \delta, q_0, F)$ un autómata finito determinista. La
función de transición $\delta$ se extiende a cadenas mediante la
\emph{función de transición extendida}
$\hat\delta : Q \times \Sigma^* \to Q$, que se define por
\[
  \hat\delta(q, \varepsilon) = q \quad\text{y}\quad
  \hat\delta(q, wa) = \delta(\hat\delta(q, w), a)
\]
para $w \in \Sigma^*$ y $a \in \Sigma$~\cite{hopcroft2006introduction}.
Así, $\hat\delta(q, u)$ es el estado en el que el autómata termina al
comenzar en $q$ y leer $u$.
\end{definicion}

Por ejemplo, en el autómata $A_0$ se tiene
\[
  \hat\delta(q_0, a) = q_1, \quad \hat\delta(q_0, aa) = q_2, \quad
  \hat\delta(q_0, aaa) = q_0 \quad\text{y}\quad \hat\delta(q_1, aa) = q_0.
\]

La función $\hat\delta$ determina, para cada cadena y cada estado de
partida, el estado en que el autómata termina, lo que permite definir el
rango sobre el autómata como un par de estados y precisar cuándo una
cadena lo realiza.

\begin{definicion}
Un \emph{rango sobre $A$} es un par de estados $(q, q') \in Q \times Q$.
Una cadena $u$ realiza el rango $(q, q')$ si $\hat\delta(q, u) = q'$.
\end{definicion}

Un rango sobre $A$ describe a la vez todas las cadenas que lo realizan.
Cuando a una variable se le asocia el rango $(q, q')$, la gramática debe
derivar la subcadena correspondiente, que a la vez debe realizar ese
rango.

En el ejemplo anterior, cada $a$ lleva al autómata $A_0$ de $q_0$ a
$q_1$, de $q_1$ a $q_2$ y de $q_2$ de vuelta a $q_0$. Así, $a$ realiza $(q_0, q_1)$,
$(q_1, q_2)$ y $(q_2, q_0)$; $aa$ realiza $(q_0, q_2)$, $(q_1, q_0)$ y
$(q_2, q_1)$; y $aaa$ realiza $(q_0, q_0)$, $(q_1, q_1)$ y $(q_2, q_2)$.

En la sección siguiente, las nociones planteadas para las cadenas se
generalizan al caso de los argumentos.

\section{El rango de un argumento sobre un autómata}\label{sec:rango-argumento}

La sección anterior estableció cuándo una cadena realiza un rango sobre
$A$. Un argumento de una cláusula es una secuencia de terminales y
variables, y las variables no son símbolos que el autómata pueda leer.
Hace falta entonces decir qué rango sobre $A$ le corresponde a un
argumento de esa forma.

\begin{definicion}
Sea $\alpha = \alpha_1 \cdots \alpha_n$ un argumento, donde cada
$\alpha_i \in T \cup V$. Una \emph{asignación de rangos sobre $A$} para
$\alpha$ es una función $\rho$ que asigna a cada $\alpha_i$ un rango
sobre $A$, $\rho(\alpha_i) = (q_i, q'_i)$, con dos condiciones:
\begin{enumerate}
  \item $q'_i = q_{i+1}$ para $1 \le i < n$;
  \item si $\alpha_i$ es un terminal $a$, entonces $\delta(q_i, a) = q'_i$.
\end{enumerate}
Cuando existe una asignación así, el argumento $\alpha$ \emph{realiza}
el rango sobre $A$ $(q_1, q'_n)$.
\end{definicion}

Por ejemplo, sobre $A_0$ se puede asignar al primer terminal $a$ del
argumento $aaX$ el rango $(q_0, q_1)$, al segundo terminal $a$ el rango
$(q_1, q_2)$, y a la variable $X$ un rango $(q_2, q')$ con $q' \in Q$.
La asignación se ve así:

\begin{center}
\begin{tikzpicture}[->, >=stealth, shorten >=1pt, auto, node distance=2.2cm, semithick]
  \node[state] (s0) {$q_0$};
  \node[state] (s1) [right=of s0] {$q_1$};
  \node[state] (s2) [right=of s1] {$q_2$};
  \node[state] (s3) [right=of s2] {$q'$};
  \path (s0) edge              node {$a$} (s1);
  \path (s1) edge              node {$a$} (s2);
  \path (s2) edge [dashed]     node {$X$} (s3);
\end{tikzpicture}
\end{center}

Las transiciones $\delta(q_0, a) = q_1$ y $\delta(q_1, a) = q_2$
satisfacen la condición 2 de la definición, que exige
$\delta(q_i, a) = q'_i$ para cada terminal $a$. El primer rango termina en $q_1$, donde empieza el segundo; el segundo
termina en $q_2$, donde empieza el rango de $X$, satisfaciendo la
condición 1. El argumento $aaX$ realiza el
rango $(q_0, q')$.

Las elecciones concretas para $X$ revelan la interacción entre la
gramática y el autómata. Si $X = \varepsilon$, entonces $aaX = aa$
realiza $(q_0, q_2)$. La asignación se ve así:

\begin{center}
\begin{tikzpicture}[->, >=stealth, shorten >=1pt, auto, node distance=1.7cm, semithick]
  \node[state] (s0) {$q_0$};
  \node[state] (s1) [right=of s0] {$q_1$};
  \node[state] (s2) [right=of s1] {$q_2$};
  \path (s0) edge node {$a$} (s1);
  \path (s1) edge node {$a$} (s2);
\end{tikzpicture}
\end{center}

$G_0$ deriva $aa$ porque tiene longitud par, pero $A_0$ no la acepta
porque no termina en el estado final $q_0$. La cadena no pertenece a la
intersección.

Si $X = aaaa$, entonces $aaX = aaaaaa$ realiza $(q_0, q_0)$. La asignación se ve así:

\begin{center}
\begin{tikzpicture}[->, >=stealth, shorten >=1pt, auto, node distance=0.9cm, semithick]
  \node[state] (s0) {$q_0$};
  \node[state] (s1) [right=of s0] {$q_1$};
  \node[state] (s2) [right=of s1] {$q_2$};
  \node[state] (s3) [right=of s2] {$q_0$};
  \node[state] (s4) [right=of s3] {$q_1$};
  \node[state] (s5) [right=of s4] {$q_2$};
  \node[state] (s6) [right=of s5] {$q_0$};
  \path (s0) edge node {$a$} (s1);
  \path (s1) edge node {$a$} (s2);
  \path (s2) edge node {$a$} (s3);
  \path (s3) edge node {$a$} (s4);
  \path (s4) edge node {$a$} (s5);
  \path (s5) edge node {$a$} (s6);
\end{tikzpicture}
\end{center}

$G_0$ deriva $aaaaaa$ y $A_0$ la acepta. La cadena pertenece a la
intersección $L(G_0) \cap L(A_0)$.

Si $X = aaa$, $G_0$ no deriva $aaa$ por su longitud impar, así que esa
elección no aparece en ninguna derivación de $G_0$.

La sección siguiente presenta la construcción de la gramática producto
sobre una subclase de gramáticas de concatenación de rango que
reconoce los mismos lenguajes.

\section{Construcción de la gramática producto}

En esta sección se desarrolla la construcción de la gramática
producto. Primero se presenta el algoritmo. A continuación se impone
una restricción sobre la gramática de entrada. Después se establece
una condición para que las ocurrencias de una misma variable reciban
el mismo rango. Luego se da el pseudocódigo. Por último se demuestra
su correctitud y se analiza su complejidad.

\subsection{Algoritmo}\label{subsec:algoritmo}

El objetivo es producir una gramática $G' = (N', T, V, P', S')$ que
reconozca las cadenas que $G$ reconoce y $A$ acepta. Una cadena $w$
pertenece a $L(A)$ si el
rango sobre $A$ que $w$ realiza parte del estado inicial y termina en
un estado final, así que $G'$ no solo debe reconocer las cadenas de
$L(G)$, sino reflejar también qué rango sobre $A$ realiza cada
argumento de cada predicado durante la derivación. Para reflejarlo,
los predicados de $G'$ se diferencian según el rango sobre $A$ que
realiza cada uno de sus argumentos: cada predicado de $G$ da lugar a
un predicado de $G'$ por cada vector de rangos sobre $A$.

\begin{definicion}
Sean $B$ un predicado de $G$ de aridad $k$ y
$\vec q = ((q_1, q'_1), \ldots, (q_k, q'_k))$ un vector de $k$ rangos
sobre $A$, uno por cada argumento de $B$. El par $(B, \vec q)$ se
llama \emph{predicado indexado} y se denota $B[\vec q]$.
\end{definicion}

Por ejemplo, como $A_0$ tiene tres estados, la gramática $G'$ de
$G_0$ sobre $A_0$ tiene nueve predicados indexados $Par[(q_i, q_j)]$,
uno por cada par $(q_i, q_j)$ de estados de $A_0$:
\[
  \begin{array}{ccc}
    Par[(q_0, q_0)], & Par[(q_0, q_1)], & Par[(q_0, q_2)], \\
    Par[(q_1, q_0)], & Par[(q_1, q_1)], & Par[(q_1, q_2)], \\
    Par[(q_2, q_0)], & Par[(q_2, q_1)], & Par[(q_2, q_2)].
  \end{array}
\]

El predicado inicial de $G'$ es un predicado $S'$ de aridad 1. El
conjunto de predicados $N'$ consta de los predicados indexados y del
predicado inicial $S'$.

Las cláusulas de $G'$ se generan a partir de las cláusulas de $G$ y de
las asignaciones de rangos sobre $A$ a sus ocurrencias de símbolos. Para
cada cláusula
\[
  B_0(\alpha_{0,1}, \ldots, \alpha_{0,k_0}) \to
  B_1(\alpha_{1,1}, \ldots, \alpha_{1,k_1}) \cdots
  B_m(\alpha_{m,1}, \ldots, \alpha_{m,k_m})
\]
de $G$ y cada asignación que cumpla las condiciones de la
sección~\ref{sec:rango-argumento}, indexando cada predicado $B_j$ por el
vector $\vec q_j$ de los rangos de sus argumentos se obtiene:
\[
  B_0[\vec q_0](\alpha_{0,1}, \ldots, \alpha_{0,k_0}) \to
  B_1[\vec q_1](\alpha_{1,1}, \ldots, \alpha_{1,k_1}) \cdots
  B_m[\vec q_m](\alpha_{m,1}, \ldots, \alpha_{m,k_m}).
\]

En la misma gramática de ejemplo, en la cláusula
$Par(aaX) \to Par(X)$ se asigna al argumento $X$ del cuerpo el rango
$(q_2, q_0)$. En la cabeza, los dos terminales llevan al autómata de
$q_0$ a $q_2$, donde empieza $X$; como $X$ realiza $(q_2, q_0)$, el
argumento $aaX$ realiza el rango $(q_0, q_0)$. La cláusula que se
obtiene en $G'$ es:
\[
  Par[(q_0, q_0)](aaX) \to Par[(q_2, q_0)](X).
\]

La cláusula $Par(aaX) \to Par(X)$ produce una cláusula en $G'$ por
cada rango asignado al argumento $X$ del cuerpo. En $A_0$, el rango de
$X$ fija el de la cabeza, así que cada una de las nueve asignaciones determina
una sola cláusula:
\[
\begin{array}{r@{\;\to\;}l}
Par[(q_0, q_0)](aaX) & Par[(q_2, q_0)](X), \\
Par[(q_0, q_1)](aaX) & Par[(q_2, q_1)](X), \\
Par[(q_0, q_2)](aaX) & Par[(q_2, q_2)](X), \\
Par[(q_1, q_0)](aaX) & Par[(q_0, q_0)](X), \\
Par[(q_1, q_1)](aaX) & Par[(q_0, q_1)](X), \\
Par[(q_1, q_2)](aaX) & Par[(q_0, q_2)](X), \\
Par[(q_2, q_0)](aaX) & Par[(q_1, q_0)](X), \\
Par[(q_2, q_1)](aaX) & Par[(q_1, q_1)](X), \\
Par[(q_2, q_2)](aaX) & Par[(q_1, q_2)](X).
\end{array}
\]
Cuando el cuerpo de una cláusula de $G$ es vacío, solo se indexa la
cabeza: la cláusula $B(\alpha_1, \ldots, \alpha_k) \to \varepsilon$
produce las cláusulas $B[\vec q](\alpha_1, \ldots, \alpha_k) \to
\varepsilon$, una por cada asignación de rangos a sus argumentos. El
argumento $\varepsilon$ realiza solo el rango $(q, q)$ para cada estado
$q$, así que la cláusula $Par(\varepsilon) \to \varepsilon$ produce
tres cláusulas en $G'$, una por cada estado del autómata:
\begin{align*}
  Par[(q_0, q_0)](\varepsilon) &\to \varepsilon, \\
  Par[(q_1, q_1)](\varepsilon) &\to \varepsilon, \\
  Par[(q_2, q_2)](\varepsilon) &\to \varepsilon.
\end{align*}

El conjunto de cláusulas $P'$ consta de las cláusulas anteriores y de
las del predicado inicial $S'$, que se presentan a continuación.

El autómata $A$ acepta una cadena cuando el rango que realiza parte
del estado inicial $q_0$ y termina en un estado final. Por eso la
construcción solo usa los predicados indexados $S[(q_0, q_F)]$ cuyo
rango parte de $q_0$ y termina en un estado final $q_F \in F$. Para
cada $q_F \in F$, $G'$ tiene una cláusula
\[
  S'(X) \to S[(q_0, q_F)](X),
\]
donde $S$ es el predicado inicial de $G$. Una cadena $w$ pertenece a
$L(G')$ si $S'(w)$ se deriva a la cadena vacía.

En el ejemplo de $G_0$ sobre $A_0$, $F = \{q_0\}$, y el rango
$(q_0, q_0)$ parte del estado inicial y termina en un estado final. El
predicado inicial $S'$ tiene la cláusula:
\[
  S'(X) \to Par[(q_0, q_0)](X).
\]

La construcción conserva todas las cláusulas que produce, tanto las
que provienen de las cláusulas de $G$ como las del predicado inicial
$S'$. Las que participan en el reconocimiento de una cadena se
determinan durante la derivación, no durante la construcción.

La construcción produce los predicados $N'$, las cláusulas $P'$ y el
predicado inicial $S'$, y mantiene los terminales $T$ y las variables
$V$ de $G$. Los cinco componentes completan la tupla
$G' = (N', T, V, P', S')$, con lo que la construcción de la gramática
producto termina.

Para ilustrar el reconocimiento en la gramática producto $G'$, se
describe paso a paso la derivación de la cadena
$aaaaaa \in L(G_0) \cap L(A_0)$:
\begin{align*}
  S'(aaaaaa) &\Rightarrow Par[(q_0, q_0)](aaaaaa) \\
          &\Rightarrow Par[(q_2, q_0)](aaaa) \\
          &\Rightarrow Par[(q_1, q_0)](aa) \\
          &\Rightarrow Par[(q_0, q_0)](\varepsilon) \\
          &\Rightarrow \varepsilon.
\end{align*}
La figura~\ref{fig:derivacion-a6} representa el rango sobre $A_0$ en
cada paso. En el primer paso se aplica la cláusula del predicado
inicial $S'$. Las cláusulas de los pasos 2, 3 y 4 provienen de
$Par(aaX) \to Par(X)$ con asignaciones distintas al argumento $aaX$, y
la del paso 5 proviene de $Par(\varepsilon) \to \varepsilon$. Como la derivación de $S'(aaaaaa)$ termina en la cadena vacía, $G'$
reconoce $aaaaaa$.

La derivación de $aaaaaa$ usa solo unas pocas de las cláusulas de
$G'$; las demás no participan en su reconocimiento. La construcción
genera, además, muchos predicados y cláusulas que no aparecen en el
reconocimiento de ninguna cadena de $L(G_0) \cap L(A_0)$. Todos los
rangos de la derivación anterior terminan en el estado final $q_0$;
los predicados indexados cuyo rango termina en otro estado, como
$Par[(q_0, q_1)]$ o $Par[(q_1, q_2)]$, quedan sin uso, junto con las
cláusulas que los tienen en la cabeza, como por ejemplo
$Par[(q_1, q_2)](aaX) \to Par[(q_0, q_2)](X)$.

\begin{figure}[H]
\centering
\setlength{\tabcolsep}{6pt}
\begin{tabular}{cc}
\begin{tikzpicture}[->, >=stealth, shorten >=1pt, auto, node distance=1.5cm, semithick, scale=0.65, transform shape]
  \node[state, initial, accepting, fill=violet!30] (q0) {$q_0$};
  \node[state] (q1) [right=of q0] {$q_1$};
  \node[state] (q2) [right=of q1] {$q_2$};
  \path (q0) edge              node {$a$} (q1)
        (q1) edge              node {$a$} (q2)
        (q2) edge [bend left=40] node {$a$} (q0);
\end{tikzpicture} &
\begin{tikzpicture}[->, >=stealth, shorten >=1pt, auto, node distance=1.5cm, semithick, scale=0.65, transform shape]
  \node[state, initial, accepting, fill=red!30] (q0) {$q_0$};
  \node[state] (q1) [right=of q0] {$q_1$};
  \node[state, fill=blue!30] (q2) [right=of q1] {$q_2$};
  \path (q0) edge              node {$a$} (q1)
        (q1) edge              node {$a$} (q2)
        (q2) edge [bend left=40] node {$a$} (q0);
\end{tikzpicture} \\[0.3em]
{\small (a) $Par[(q_0, q_0)](aaaaaa)$} &
{\small (b) $Par[(q_2, q_0)](aaaa)$} \\[1.2em]
\begin{tikzpicture}[->, >=stealth, shorten >=1pt, auto, node distance=1.5cm, semithick, scale=0.65, transform shape]
  \node[state, initial, accepting, fill=red!30] (q0) {$q_0$};
  \node[state, fill=blue!30] (q1) [right=of q0] {$q_1$};
  \node[state] (q2) [right=of q1] {$q_2$};
  \path (q0) edge              node {$a$} (q1)
        (q1) edge              node {$a$} (q2)
        (q2) edge [bend left=40] node {$a$} (q0);
\end{tikzpicture} &
\begin{tikzpicture}[->, >=stealth, shorten >=1pt, auto, node distance=1.5cm, semithick, scale=0.65, transform shape]
  \node[state, initial, accepting, fill=violet!30] (q0) {$q_0$};
  \node[state] (q1) [right=of q0] {$q_1$};
  \node[state] (q2) [right=of q1] {$q_2$};
  \path (q0) edge              node {$a$} (q1)
        (q1) edge              node {$a$} (q2)
        (q2) edge [bend left=40] node {$a$} (q0);
\end{tikzpicture} \\[0.3em]
{\small (c) $Par[(q_1, q_0)](aa)$} &
{\small (d) $Par[(q_0, q_0)](\varepsilon)$}
\end{tabular}
\caption{Rango sobre $A_0$ en cada paso de la derivación de $aaaaaa$.
El color azul corresponde al estado de partida del rango; el rojo,
al de llegada; el violeta, a ambos cuando coinciden.}
\label{fig:derivacion-a6}
\end{figure}

Por conveniencia, la construcción anterior se restringe a una clase
de gramáticas de entrada que evita casos aparte. La sección siguiente
la precisa.

\subsection{Restricción sobre la gramática de entrada}

Cada predicado se indexa por los rangos sobre $A$ de sus argumentos.
El cuerpo de una cláusula fija los rangos de las variables que
contiene. Si una variable apareciera en la cabeza pero no en el
cuerpo, el cuerpo no fijaría su rango, y habría que tratarla como un
caso aparte.

Para evitarlo, las gramáticas de entrada se restringen a las que
cumplen una propiedad sobre sus cláusulas: en cada cláusula, toda
variable de la cabeza aparece también en el cuerpo.

\begin{definicion}
Una gramática de concatenación de rango es \emph{top-down non-erasing}
si en cada cláusula toda variable de la cabeza aparece también en el
cuerpo.
\end{definicion}

Se pide además que los argumentos del cuerpo sean variables, de modo
que cada variable del cuerpo sea un argumento por sí sola.

\begin{definicion}
Una gramática de concatenación de rango es \emph{non-combinatorial} si
en cada cláusula los argumentos del cuerpo son variables.
\end{definicion}

Por ejemplo, las cláusulas de $G_0$ cumplen ambas: $Par(aaX) \to
Par(X)$ tiene la variable $X$ en la cabeza y en el cuerpo, donde es un
argumento por sí sola, y $Par(\varepsilon) \to \varepsilon$ no tiene
variables. Una cláusula como $A(XY) \to B(X)$ no es \emph{top-down
non-erasing}, pues $Y$ figura en la cabeza y falta en el cuerpo, y
$A(X) \to B(Xa)$ no es \emph{non-combinatorial}, pues el argumento del
cuerpo $Xa$ no es una variable.

Exigir ambas propiedades no descarta ningún lenguaje: toda gramática de
concatenación de rango admite una equivalente que las
cumple~\cite{boullier1998proposal, kallmeyer2010parsing}.

Además de la restricción sobre la gramática de entrada, hace falta una
condición sobre los rangos que se asignan. La sección siguiente la
precisa.

\subsection{Variables repetidas en una cláusula}\label{subsec:coherencia}

En las gramáticas de concatenación de rango, una misma variable puede
aparecer varias veces en una cláusula. En la construcción de la
sección~\ref{subsec:algoritmo}, cada argumento del cuerpo se indexa por
separado,
sin que se fije el mismo rango para las ocurrencias de una misma
variable, que representan una sola cadena.
Sin una condición que lo imponga, se les pueden asignar rangos
distintos, y entonces $G'$ reconocería cadenas que no están en la
intersección.

Por eso es necesario añadir la siguiente condición: en cada cláusula,
a todas las ocurrencias de una variable se les asigna el mismo rango.
Más formalmente, si una variable $X$ aparece en dos argumentos del
cuerpo con asignaciones de rangos $\rho$ y $\rho'$, entonces
$\rho(X) = \rho'(X)$.

Por ejemplo, si a $G_0$ se le agrega un predicado $Cinco$, que reconoce
las cadenas de longitud múltiplo de cinco, y una cláusula inicial que
combina $Par$ y $Cinco$ sobre la misma variable, la gramática
resultante es $G_c = (N, T, V, P, S)$, donde $N = \{S, Par, Cinco\}$,
$T = \{a\}$, $V = \{X\}$ y las cláusulas de $P$ son
\[
\begin{aligned}
  S(X) &\to Par(X)\,Cinco(X), \\
  Par(aaX) &\to Par(X), \quad Par(\varepsilon) \to \varepsilon, \\
  Cinco(aaaaaX) &\to Cinco(X), \quad Cinco(\varepsilon) \to \varepsilon.
\end{aligned}
\]
$G_c$ reconoce las cadenas que $Par$ y $Cinco$ reconocen a la vez, las
de longitud múltiplo de diez.

Sobre $A_0$, la cadena $a^{10}$ realiza el rango $(q_0, q_1)$, porque su
longitud deja resto uno al dividir entre tres, y el rango no termina en
el estado final. Su longitud es múltiplo de dos y de cinco, pero no de
tres, así que no pertenece a la intersección.

Sin añadir la condición existe la cláusula
\[
  S[(q_0, q_0)](X) \to Par[(q_0, q_1)](X)\,Cinco[(q_0, q_1)](X),
\]
con el predicado de la cabeza indexado en el rango $(q_0, q_0)$ y los
predicados del cuerpo indexados en el rango $(q_0, q_1)$.
Como $a^{10}$ realiza $(q_0, q_1)$ y la reconocen tanto $Par$ como
$Cinco$, los dos cuerpos se derivan en la cadena vacía, y $G'$ reconoce
$a^{10}$, la cual está fuera de la intersección.

Con la condición, las ocurrencias de $X$ llevan el mismo rango, y la
única cláusula válida con la cabeza indexada en el rango $(q_0, q_0)$ es
\[
  S[(q_0, q_0)](X) \to Par[(q_0, q_0)](X)\,Cinco[(q_0, q_0)](X).
\]
Como
$a^{10}$ realiza $(q_0, q_1)$ y no $(q_0, q_0)$,
$Par[(q_0, q_0)](a^{10})$ no se deriva y $a^{10}$ queda fuera del
lenguaje de $G'$, como debería ser. La gramática producto reconoce
entonces exactamente las cadenas de longitud múltiplo de dos, tres y
cinco, es decir de treinta.

La repetición de variables es justo lo que separa a las gramáticas de
concatenación de rango generales de las simples, y la condición no
hace falta para las simples: en una gramática simple
cada variable aparece una sola vez en el cuerpo, así que no hay varias
ocurrencias que coordinar. Imponerla es lo que permite tratar
gramáticas de concatenación de rango generales.

\subsection{Pseudocódigo}

Ya se presentaron la idea del algoritmo, la restricción sobre la
gramática de entrada y la condición sobre las ocurrencias de una misma
variable; con ellas se fija la construcción. A continuación se describe
el pseudocódigo.

Recibe una gramática de concatenación de rango non-combinatorial y
top-down non-erasing $G = (N, T, V, P, S)$ y un autómata finito determinista
$A = (Q, \Sigma, \delta, q_0, F)$; devuelve $G' = (N', T, V, P', S')$
con $L(G') = L(G) \cap L(A)$.

\begin{codebox}
\Procname{$\proc{Product-Grammar}(G, A)$}
\li $P' = \emptyset$
\li \For each clause $c = B_0(\alpha_0) \to B_1(\alpha_1) \cdots B_m(\alpha_m)$ in $P$
\li \Do
        let $\sigma_1, \sigma_2, \ldots$ be the symbol occurrences of $c$ and its empty arguments
\li     $\ell = $ number of these
\li     \For each $\mu = ((q_1, q'_1), \ldots, (q_\ell, q'_\ell)) \in (Q \times Q)^\ell$
\li     \Do
            \If $\proc{Consistent}(c, \mu)$
\li         \Then
                \For $j = 0$ \To $m$
\li             \Do
                    $\vec{q}_j = (\proc{Range}(\alpha_{j,1}, \mu), \ldots, \proc{Range}(\alpha_{j,k_j}, \mu))$
                \End
\li             add $B_0[\vec{q}_0](\alpha_0) \to B_1[\vec{q}_1](\alpha_1) \cdots B_m[\vec{q}_m](\alpha_m)$ to $P'$
            \End
        \End
    \End
\li \For each $q \in F$
\li \Do
        add $S'(X) \to S[(q_0, q)](X)$ to $P'$
    \End
\li $N' = \{\, S' \,\} \cup \{\, B' : B' \text{ occurs in } P' \,\}$
\li \Return $(N', T, V, P', S')$
\end{codebox}

En la línea~2 se recorren las cláusulas de $G$. Para cada una, en las
líneas~3 y~4 se cuentan los símbolos que aparecen en ella, con
repeticiones, y sus argumentos vacíos. En la línea~5 se recorren todas
las posibles asignaciones de rangos. De la línea~6 a la~9, con $\proc{Consistent}$ se conservan
solo las asignaciones que cumplen las condiciones de las
secciones~\ref{sec:rango-argumento} y~\ref{subsec:coherencia}; por
cada una se indexa cada predicado por los rangos de sus argumentos con
$\proc{Range}$, y se agrega la cláusula a $P'$. En las líneas~10 y~11 se
agregan las cláusulas del predicado inicial $S'$, una por cada estado
final. En la línea~12 se agregan a $N'$ el
predicado $S'$ y los predicados que aparecen en $P'$.

En $\proc{Consistent}$ se comprueban las condiciones sobre $\mu$:

\begin{codebox}
\Procname{$\proc{Consistent}(c, \mu)$}
\li \For each non-empty argument $\alpha = \alpha_1 \cdots \alpha_n$ of $c$, with $\mu(\alpha_i) = (q_i, q'_i)$
\li \Do
        \For $i = 1$ \To $n$
\li     \Do
            \If $i < n$ \kw{and} $q'_i \neq q_{i+1}$
\li         \Then \Return \const{false}
            \End
\li         \If $\alpha_i$ is a terminal $a$ \kw{and} $\delta(q_i, a) \neq q'_i$
\li         \Then \Return \const{false}
            \End
        \End
    \End
\li \For each empty argument $\alpha$ of $c$, with $\mu(\alpha) = (q, q')$
\li \Do
        \If $q \neq q'$
\li     \Then \Return \const{false}
        \End
    \End
\li \For each variable $X$ of $c$
\li \Do
        \If two occurrences of $X$ have different ranges under $\mu$
\li     \Then \Return \const{false}
        \End
    \End
\li \Return \const{true}
\end{codebox}

Las líneas~1--9 de $\proc{Consistent}$ son la condición de la
sección~\ref{sec:rango-argumento}; las líneas~10--12, la de la
sección~\ref{subsec:coherencia}.

$\proc{Range}(\alpha, \mu)$ es el rango que el argumento $\alpha$
realiza bajo $\mu$:

\begin{codebox}
\Procname{$\proc{Range}(\alpha, \mu)$}
\li \If $\alpha = \varepsilon$
\li \Then \Return $\mu(\alpha)$
    \End
\li let $\alpha = \alpha_1 \cdots \alpha_n$, with $\mu(\alpha_i) = (q_i, q'_i)$
\li \Return $(q_1, q'_n)$
\end{codebox}

Para $\alpha = \alpha_1 \cdots \alpha_n$ es $(q_1, q'_n)$, el inicio del
primer símbolo y el final del último; para $\alpha = \varepsilon$ es el
rango $(q, q)$ que $\mu$ le asigna, el rango que $\alpha$ realiza según
la sección~\ref{sec:rango-argumento}.

\subsection{Correctitud}

Las secciones anteriores describen la construcción de la gramática
producto $G'$. Falta demostrar que $G'$ reconoce exactamente
$L(G) \cap L(A)$. En toda la sección $G$ es non-combinatorial y top-down
non-erasing, y $G'$ es la gramática que la construcción produce a partir
de $G$ y $A$.

La igualdad $L(G') = L(G) \cap L(A)$ se prueba por las dos inclusiones.
La clave es una propiedad de los índices: una cláusula indexada solo
reconoce cadenas que realizan los rangos con que se indexan sus
argumentos. De ahí, $A$ acepta toda cadena que $G'$ reconoce; y, al
borrar los índices, también la reconoce $G$. En el otro sentido, se
parte de una derivación de $G$ y se indexa cada predicado con los rangos
sobre $A$ que sus argumentos realizan al leer la cadena; el resultado es
una derivación de $G'$.

Sea $B[\vec q](u_1, \ldots, u_k)$ el predicado indexado $B[\vec q]$
aplicado a las cadenas $u_1, \ldots, u_k$, una por cada argumento, con
$\vec q = ((q_1, q'_1), \ldots, (q_k, q'_k))$.

\begin{lema}\label{lema:proyeccion}
Si $B[\vec q](u_1, \ldots, u_k)$ se deriva a $\varepsilon$ en $G'$,
entonces $B(u_1, \ldots, u_k)$ se deriva a $\varepsilon$ en $G$.
\end{lema}

\begin{proof}
Las cláusulas que intervienen en una derivación de
$B[\vec q](u_1, \ldots, u_k)$ son cláusulas indexadas, cada una
proveniente de una cláusula de $P$ a la que solo se le agregan los
índices de los predicados; el predicado $S'$ no aparece, pues no figura
en el cuerpo de ninguna cláusula. Al borrar los índices, cada cláusula
indexada se reduce a su cláusula de $P$, y se obtiene una derivación de
$B(u_1, \ldots, u_k)$ en $G$.
\end{proof}

La hipótesis non-combinatorial es clave en el segundo lema.

\begin{lema}\label{lema:realizacion}
Si $B[\vec q](u_1, \ldots, u_k)$ se deriva a $\varepsilon$ en $G'$,
entonces cada $u_l$ realiza el rango $(q_l, q'_l)$.
\end{lema}

\begin{proof}
La prueba es por inducción en la cantidad de pasos de la derivación. La
derivación comienza con una cláusula de $P'$ de cabeza $B[\vec q]$, que
proviene de una cláusula
$c = B(\alpha_0) \to B_1(\alpha_1) \cdots B_m(\alpha_m)$ de $P$ y de una
asignación $\mu$ con $\proc{Consistent}(c, \mu)$:
\[
  B[\vec q](\alpha_0) \to B_1[\vec q_1](\alpha_1) \cdots B_m[\vec q_m](\alpha_m),
\]
con $(q_l, q'_l) = \proc{Range}(\alpha_{0,l}, \mu)$. La cláusula se instancia
con una sustitución $\theta$ de las variables por subcadenas, de modo
que $u_l = \theta(\alpha_{0,l})$ y, en el cuerpo, los predicados
$B_j[\vec q_j](\theta(\alpha_{j,1}), \ldots, \theta(\alpha_{j,k_j}))$ se
derivan a $\varepsilon$ en menos pasos.

Como $G$ es non-combinatorial, cada argumento del cuerpo es una variable
$X$, cuyo rango $\mu(X)$ es el mismo en todas sus ocurrencias
(sección~\ref{subsec:coherencia}) y cumple $\proc{Range}(X, \mu) =
\mu(X)$. Por hipótesis de inducción, $\theta(X)$ realiza $\mu(X)$ para
toda variable $X$ del cuerpo; como $G$ es top-down non-erasing, lo mismo
vale para toda variable de $c$.

Sea $\alpha_{0,l}$ un argumento de la cabeza. Si $\alpha_{0,l} =
\varepsilon$, entonces $(q_l, q'_l) = \proc{Range}(\varepsilon, \mu) =
\mu(\varepsilon) = (p, p)$ por $\proc{Consistent}$, y $u_l = \varepsilon$
realiza $(p, p)$. Si $\alpha_{0,l} = \beta_1 \cdots \beta_s$ con
$\mu(\beta_t) = (p_t, p'_t)$, entonces cada $\theta(\beta_t)$ realiza
$(p_t, p'_t)$: si $\beta_t$ es un terminal $a$, por la condición
$\delta(p_t, a) = p'_t$ de $\proc{Consistent}$; si $\beta_t$ es una
variable $X$, entonces $(p_t, p'_t) = \mu(X)$ por la coherencia
(sección~\ref{subsec:coherencia}) y $\theta(\beta_t) = \theta(X)$
realiza $\mu(X)$. Como $p'_t = p_{t+1}$ para $1 \le t < s$ por
$\proc{Consistent}$, la concatenación $u_l = \theta(\beta_1) \cdots
\theta(\beta_s)$ realiza $(p_1, p'_s) = \proc{Range}(\alpha_{0,l}, \mu) =
(q_l, q'_l)$.
\end{proof}

Con los dos lemas se obtiene el teorema.

\begin{teorema}\label{teo:correctitud}
$L(G') = L(G) \cap L(A)$.
\end{teorema}

\begin{proof}
\emph{Inclusión $L(G') \subseteq L(G) \cap L(A)$.} Sea $w \in L(G')$;
entonces $S'(w)$ se deriva a $\varepsilon$. Las únicas cláusulas de
cabeza $S'$ son $S'(X) \to S[(q_0, q_F)](X)$ con $q_F \in F$, así que el
primer paso es $S'(w) \to S[(q_0, q_F)](w)$ para algún $q_F \in F$, y
$S[(q_0, q_F)](w)$ se deriva a $\varepsilon$. Por el
lema~\ref{lema:proyeccion}, $S(w)$ se deriva a $\varepsilon$ en $G$, esto
es, $w \in L(G)$. Por el lema~\ref{lema:realizacion}, $w$ realiza
$(q_0, q_F)$, es decir, $\hat\delta(q_0, w) = q_F$; como $q_F \in F$,
$w \in L(A)$. Luego $w \in L(G) \cap L(A)$.

\emph{Inclusión $L(G) \cap L(A) \subseteq L(G')$.} Sea
$w = a_1 \cdots a_n \in L(G) \cap L(A)$. En la corrida de $A$ sobre $w$,
los estados $\hat g_0, \ldots, \hat g_n$ cumplen $\hat g_0 = q_0$ y
$\hat g_i = \delta(\hat g_{i-1}, a_i)$; como $w \in L(A)$,
$\hat g_n \in F$. La subcadena que ocupa las posiciones $i$ a $j$ realiza
el rango sobre $A$ $(\hat g_i, \hat g_j)$.

Como $w \in L(G)$, hay una derivación de $S(w)$ a $\varepsilon$ en $G$,
en la que cada predicado se aplica a subcadenas de $w$ en posiciones
fijas. Al indexar cada predicado por los rangos sobre $A$ de sus
argumentos se obtiene una derivación en $G'$, porque cada cláusula
indexada pertenece a $P'$. En efecto, sea $B_0(\alpha_0) \to
B_1(\alpha_1) \cdots B_m(\alpha_m)$ una cláusula instanciada en la
derivación, y $\mu$ la asignación que a cada ocurrencia de símbolo le
asigna el rango sobre $A$ de la posición que ocupa. La asignación $\mu$
cumple $\proc{Consistent}$: símbolos contiguos de un argumento ocupan
posiciones contiguas, así que el final de un rango coincide con el
inicio del siguiente; para un terminal $a$ en la posición $i{+}1$ se
tiene $\delta(\hat g_i, a) = \hat g_{i+1}$; un argumento vacío en la
posición $i$ tiene el rango $(\hat g_i, \hat g_i)$; y las ocurrencias de
una variable ocupan las mismas posiciones, luego les corresponde el
mismo rango sobre $A$. Por tanto, la construcción produce la cláusula
indexada, que pertenece a $P'$.

La cabeza $S(w)$ se indexa por el rango $(\hat g_0, \hat g_n) =
(q_0, \hat g_n)$. Como $\hat g_n \in F$, la cláusula
$S'(X) \to S[(q_0, \hat g_n)](X)$ pertenece a $P'$, así que $S'(w)$ se
deriva a $\varepsilon$ y $w \in L(G')$.
\end{proof}

\bibliographystyle{plain}
\bibliography{../bib/referencias}

\end{document}
